% Question 2 formulas for paper handoff

\section{问题二模型公式}

\subsection{决策变量}

设
\[
m_1,m_2\in\{0,1,2\},\qquad y,z\in\{0,1\}.
\]

其中 \(m_1,m_2\) 表示零配件 1、零配件 2 的检测策略，\(y\) 表示是否检测成品，\(z\) 表示是否拆解不合格成品。

零配件策略定义为
\[
m_i=
\begin{cases}
0, & \text{第一次检测，拆解后不检测},\\
1, & \text{第一次不检测，拆解后也不检测},\\
2, & \text{第一次不检测，拆解后检测}.
\end{cases}
\]

\subsection{目标函数}

以最终交付 1 件合格成品为核算单位，设市场售价为 \(S\)，期望成本为 \(E\)，则期望利润为
\[
\Pi=S-E.
\]

最优策略为
\[
S^\ast=\arg\max_{S_i}\Pi(S_i),
\]
或等价地
\[
S^\ast=\arg\min_{S_i}E(S_i).
\]

\subsection{零配件检测策略}

先枚举 9 种零配件策略：
\[
(m_1,m_2),\qquad m_1,m_2\in\{0,1,2\}.
\]

若零配件 \(i\) 第一次不检测，则进入装配环节的期望成本和合格概率分别为
\[
K_i=c_i,\qquad u_i=1-p_i.
\]

若零配件 \(i\) 第一次检测，则不合格件被剔除，需要重复购买和检测直到获得合格件。设购买单价为 \(c_i\)，检测成本为 \(t_i\)，次品率为 \(p_i\)，则
\[
K_i
=
\sum_{k=1}^{\infty}k(c_i+t_i)p_i^{k-1}(1-p_i)
=
\frac{c_i+t_i}{1-p_i},
\qquad
u_i=1.
\]

\subsection{成品不合格概率}

设成品装配次品率为 \(p_f\)。在给定 \((m_1,m_2)\) 后，首次装配时成品合格概率为
\[
g=u_1u_2(1-p_f).
\]

成品不合格概率为
\[
q=1-g.
\]

\subsection{成品检测判断}

成品检测成本为 \(t_f\)，不合格品流入市场造成的调换损失为 \(L\)。若不检测成品，单件成品带来的期望调换损失为 \(qL\)；若检测成品，需要支付检测成本 \(t_f\)。

因此成品检测的局部占优条件为
\[
t_f<qL.
\]

若该式成立，取
\[
y=1;
\]
否则取
\[
y=0.
\]

\subsection{拆解后的零配件处理}

拆解不会损坏零配件，但拆解后的零配件是否已知合格取决于首次是否检测。

若 \(m_i=0\)，则零配件 \(i\) 首次检测后才进入装配，因此拆解后仍为已知合格件，不需要重复检测。

若 \(m_i=1\)，则零配件 \(i\) 首次不检测，拆解后也不检测，直接进入下一轮装配。

若 \(m_i=2\)，则零配件 \(i\) 首次不检测，但拆解后检测。若检测出不合格，则丢弃并重新购买；若检测合格，则进入下一轮装配。

\subsection{期望成本递推}

为了处理拆解后的零配件状态，设
\[
E(s_1,s_2)
\]
表示在零配件 1 状态为 \(s_1\)、零配件 2 状态为 \(s_2\) 时，最终交付 1 件合格成品所需的期望成本。

零配件状态包括：
\[
s_i\in\{\text{无可用件},\text{已知合格},\text{未知但实际合格},\text{未知且实际不合格}\}.
\]

对每个状态，按照该策略下的购买、检测、装配、成品检测、调换、拆解流程列出递推方程：
\[
E(s_1,s_2)
=
C(s_1,s_2)
+
\sum_{s_1',s_2'}P(s_1',s_2'\mid s_1,s_2)E(s_1',s_2').
\]

将所有状态的递推方程联立，得到线性方程组，求解初始状态 \(E(\text{无可用件},\text{无可用件})\)。

期望利润为
\[
\Pi=S-E(\text{无可用件},\text{无可用件}).
\]

\subsection{降维求解流程}

综上，求解流程为
\[
(m_1,m_2)\rightarrow q\rightarrow y\rightarrow z\rightarrow \Pi.
\]

局部判断可以减少枚举数量，但最终最优策略仍需要通过总期望利润比较：
\[
\boxed{\text{局部判断可以减少枚举数量，但最终最优策略仍要靠总期望利润比较。}}
\]
